\section{Scope and open problems}\label{sec:boundary}

The recursive carrier theorem closes the reduced contained-component problem
at arbitrary redundancy.  The remaining obstruction is arithmetic:
small-characteristic points on later maximal Lucas carriers need not share
the degree-nine answer.
\begin{enumerate}[leftmargin=2em]
\item Redundancy five is complete for every $q\geq7$.
\item Redundancy six is a complete deep-hole classification over every
      prime power \(q\geq7\).  Redundancy seven has a complete split-free
      classification over the same range; at $q=8$ its radius is seven and
      exact extraction gives ten deep directions in the diagonal tangent and
      central-nucleus orbits; only the $q=7,9$ radius gates remain separate.
\item Redundancies eight, nine, and ten have exact persistent-only
      classifications in the stated high-field ranges.
\item At arbitrary redundancy, the reduced carrier and its degree-one modular
      union are exact.  Split-squarefree arithmetic on later Lucas carriers
      remains open outside the evaluated levels.
\end{enumerate}

In particular, the paper does not settle the general Reed--Solomon
deep-hole conjecture.  It proves the reduced carrier and the numerical
split-free containment unconditionally in the displayed field range.  When
the characteristic exceeds \(r-1\), the Lucas carrier is empty and the
separate coding-theoretic radius gate promotes that containment to the exact
high-weight-coset classification.

The last sporadic fields at redundancies five, six, and seven are
\(19,13,11\), while the corresponding first prime-power field-order gates
are \(23,29,37\).  These opposite trends ask whether sufficiently high
redundancy has any sporadic fields outside the persistent and modular
loci.
